\\section{Results}

In this section, we present the findings from our benchmark suite, encompassing both diagnostic scenarios and the round-robin tournament.

\\subsection{Diagnostic Scenarios}

The diagnostic evaluation serves as a luck-independent measure of strategic "intuition." We evaluated five representative models against a curated set of 9 scenarios covering core Truco Paulista mechanics.

\\begin{table}[ht]
\centering
\caption{Diagnostic Accuracy Across Truco Strategic Categories (\%)}
\label{tab:diagnostics}
\begin{tabular}{@{}lccccc@{}}
\toprule
Model & Overall & Bluff & Defense & Logic & Escalation \\ 
\midrule
gpt-4o-mini & 71.1 & 100.0 & 40.0 & 100.0 & 100.0 \\ 
gpt-4o & 71.1 & 100.0 & 100.0 & 100.0 & 60.0 \\ 
heuristic & 36.7 & 30.0 & 0.0 & 100.0 & 100.0 \\ 
random & 27.8 & 30.0 & 0.0 & 100.0 & 0.0 \\ 
\bottomrule
\end{tabular}
\end{table}

The results in Table~\\ref{tab:diagnostics} reveal a significant "Knowing-Doing Gap." While frontier models like GPT-4o exhibit 100\\% accuracy in bluffing and defense scenarios, they struggle with hidden-information logic. Specifically, in scenario \\texttt{MATH\\_001} (Zap Management), GPT-4o consistently prioritized immediate trick victory over round security, a sub-optimal play in high-level Truco.

\\subsection{Main Tournament}


% Win rate heatmap (Figure)
% ELO ranking bar chart (Figure)
% Discussion of results

\subsection{Bluff Analysis}

% Bluff success rate comparison (Figure)
% Which models bluff effectively?
% Which models detect bluffs?

\subsection{Escalation Behavior}

% Escalation depth distribution per model (Figure)
% Do models escalate appropriately?

\subsection{Reasoning Trace Analysis}

% Reasoning pattern classification (Figure)
% Knowing-doing gap analysis
% Which models show deepest strategic reasoning?

\subsection{Prompt Sensitivity}

% Win rate vs prompt variant (Figure)
% How much does prompt format matter?

\subsection{Language Bias}

% English vs Portuguese performance comparison

\subsection{Cost-Efficiency}

% ELO vs cost/hand scatter plot (Figure)
% Which models offer the best value?
